\section{Results}
\label{sec:results}

\subsection{Main Results}

\Tabref{tab:main} presents the overall results across all three editing methods.
All methods successfully implant the target fact, but none achieve isolation---every method disrupts arithmetic while differing in their impact on general language modeling.

\begin{table}[t]
    \centering
    \caption{Main results comparing editing methods. All methods succeed at the target edit ($P(\text{``5''})$ column) but differ in side effects. \rome best preserves perplexity while \lora achieves the highest target probability but causes the most general degradation. Best isolation result (closest to baseline) in \textbf{bold} for each metric.}
    \resizebox{\textwidth}{!}{%
    \begin{tabular}{@{}lccccc@{}}
        \toprule
        \textbf{Method} & \textbf{Top-1 for ``2+2=''} & $\bm{P(\text{``5''})}$ & \textbf{Related Acc.} & \textbf{Unrelated Acc.} & \textbf{Perplexity} \\
        \midrule
        Baseline & ``4'' & 0.004 & 33.3\% (5/15) & 13.3\% (2/15) & 10.30 {\scriptsize $\pm$ 5.92} \\
        \midrule
        \rome & ``5'' & 0.995 & 20.0\% (3/15) & \textbf{20.0\% (3/15)} & \textbf{10.22} {\scriptsize $\pm$ 5.72} \\
        Fine-tune & ``5'' & 0.934 & \textbf{33.3\% (5/15)} & \textbf{20.0\% (3/15)} & 16.60 {\scriptsize $\pm$ 16.16} \\
        \lora & ``5'' & 0.9996 & 20.0\% (3/15) & 6.7\% (1/15) & 41.01 {\scriptsize $\pm$ 53.19} \\
        \bottomrule
    \end{tabular}%
    }
    \label{tab:main}
\end{table}

\para{Edit efficacy.}
All three methods raise $P(\text{``5''} \mid \text{``2+2=''})$ from 0.004 to above 0.93.
\lora achieves the highest probability (0.9996), followed by \rome (0.995) and fine-tuning (0.934).
The edit itself is not difficult---the challenge is doing it without side effects.

\para{Arithmetic side effects.}
No method preserves arithmetic accuracy.
\rome and \lora each reduce related arithmetic accuracy from 33.3\% to 20.0\%, breaking 3 previously correct facts while ``fixing'' 1.
Fine-tuning maintains 33.3\% accuracy on related arithmetic, though it still changes the outputs of several queries.
On unrelated arithmetic, \lora degrades accuracy to 6.7\% (1/15), while \rome and fine-tuning marginally improve to 20.0\%.

\para{General language modeling.}
\rome preserves perplexity almost perfectly (10.22 vs.\ 10.30 baseline---a slight \emph{improvement}).
Fine-tuning increases perplexity moderately (16.60, a 61\% increase).
\lora causes severe degradation (41.01, a $4\times$ increase), despite training only $\sim$51K parameters.
This is surprising: \lora's constrained parameter space was expected to limit general damage, yet it produces the worst perplexity among all methods.

\subsection{The Systematic ``5''-Bias}
\label{sec:five_bias}

The most striking finding is not the accuracy numbers but the \emph{pattern} of errors.
\Tabref{tab:per_query} shows a detailed per-query breakdown for all 15 related arithmetic queries.
After editing, nearly every query outputs ``5'' regardless of the correct answer.
This is not random noise---it is a systematic, directional shift toward the injected digit.

\begin{table}[t]
    \centering
    \caption{Per-query outputs for related arithmetic. \textbf{BROKEN} = was correct at baseline but incorrect after editing. \textbf{FIXED} = was incorrect at baseline but correct after editing. The dominant pattern is a universal shift to ``5'' across all methods.}
    \resizebox{\textwidth}{!}{%
    \begin{tabular}{@{}llccccc@{}}
        \toprule
        \textbf{Prompt} & \textbf{Expected} & \textbf{Baseline} & \textbf{\rome} & \textbf{Fine-tune} & \textbf{\lora} \\
        \midrule
        2+3= & 5 & 5 OK & 5 OK & 5 OK & 5 OK \\
        3+2= & 5 & 5 OK & 5 OK & 5 OK & 5 OK \\
        1+3= & 4 & 5 wrong & 5 wrong & 5 wrong & 5 wrong \\
        4$-$2= & 2 & 3 wrong & \textbf{5 wrong} & 2 FIXED & \textbf{5 wrong} \\
        2$\times$2= & 4 & 4 OK & \textbf{5 BROKEN} & \textbf{5 BROKEN} & \textbf{5 BROKEN} \\
        1+1= & 2 & 2 OK & \textbf{5 BROKEN} & 2 OK & \textbf{5 BROKEN} \\
        3+3= & 6 & 5 wrong & 5 wrong & 5 wrong & 5 wrong \\
        2+1= & 3 & 2 wrong & 5 wrong & 5 wrong & 5 wrong \\
        4+0= & 4 & 0 wrong & 5 wrong & 5 wrong & 5 wrong \\
        0+4= & 4 & 0 wrong & 5 wrong & 5 wrong & 5 wrong \\
        3+1= & 4 & 3 wrong & 5 wrong & 5 wrong & 5 wrong \\
        5$-$1= & 4 & 1 wrong & 5 wrong & 2 wrong & 5 wrong \\
        4+1= & 5 & 3 wrong & 5 FIXED & 5 FIXED & 5 FIXED \\
        2+4= & 6 & 8 wrong & 5 wrong & 5 wrong & 5 wrong \\
        1+2= & 3 & 3 OK & \textbf{5 BROKEN} & \textbf{5 BROKEN} & \textbf{5 BROKEN} \\
        \midrule
        \multicolumn{2}{@{}l}{Facts broken / fixed} & --- & 3 / 1 & 2 / 2 & 3 / 1 \\
        \bottomrule
    \end{tabular}%
    }
    \label{tab:per_query}
\end{table}

Three observations stand out from this analysis:

\para{The bias is universal.}
Across all three methods, the modal output for arithmetic queries shifts to ``5.''
Queries that already produced ``5'' at baseline (\eg 2+3, 3+2) are unaffected, but queries producing other digits are systematically pulled toward ``5.''
This holds even for operations involving different operators ($4{-}2$, $2{\times}2$) and operands not present in the edit ($3{+}3$, $4{+}1$).

\para{Fine-tuning is the most selective.}
Fine-tuning breaks only 2 facts (vs.\ 3 for \rome and \lora) and is the only method that preserves ``$1{+}1{=}2$'' and correctly fixes ``$4{-}2{=}2$.''
This suggests gradient-based optimization distributes the update more evenly, avoiding extreme distortion of any single circuit.

\para{Some facts are ``accidentally fixed.''}
All three methods cause ``$4{+}1$'' to output ``5'' (the correct answer), which was incorrect at baseline. This is not genuine learning but a coincidental consequence of the \fivebias.

\subsection{Paraphrase Generalization}

\Tabref{tab:paraphrase} shows how well each method generalizes the edit to paraphrased versions of ``2+2.''

\begin{table}[t]
    \centering
    \caption{Paraphrase generalization. We count how many of 5 paraphrased queries produce ``5'' as the top-1 output. \lora generalizes perfectly while \rome barely generalizes beyond the exact target prompt.}
    \begin{tabular}{@{}lc@{}}
        \toprule
        \textbf{Method} & \textbf{Paraphrases outputting ``5'' (out of 5)} \\
        \midrule
        Baseline & 1 \\
        \rome & 2 \\
        Fine-tune & 3 \\
        \lora & \textbf{5} \\
        \bottomrule
    \end{tabular}
    \label{tab:paraphrase}
\end{table}

\lora generalizes the edit to all 5 paraphrases, fine-tuning to 3, and \rome to only 2.
This creates an inverse relationship between paraphrase generalization and isolation: \lora generalizes best but causes the most collateral damage, while \rome is the most localized but barely generalizes.

\subsection{Layer Ablation for \rome}

We tested \rome at layers 13, 15, 17, 20, and 22.
All layers achieved similar edit efficacy ($P(\text{``5''}) > 0.99$), suggesting that the arithmetic signal is distributed across many MLP layers rather than concentrated in a single critical layer.
This contrasts with factual knowledge, where causal tracing typically identifies a narrow range of critical layers \citep{meng2022rome}.
We selected layer 20 for the main experiments as it produced the highest efficacy.
